\documentclass{book}
\usepackage[paperwidth=5.5in,paperheight=8.5in]{geometry}
\usepackage{lettrine} % for Drop Caps
\usepackage{GoudyIn}
\renewcommand{\LettrineFontHook}{\GoudyInfamily{}}
% NOTE: use PDFMixTool - Page Layout, to nicely format this to a 2-page spread
\setlength{\parskip}{1em}

\begin{document}

\title{A Journeyman's Trade\\
\vspace{0.5em} \large an excerpt from \textit{Ireland: A Novel}}
\author{Frank Delaney}
\date{}%blank for no date
\maketitle


\begin{center}
    \LARGE \textbf{A Journeyman's Trade} \\
\end{center}
\vspace{0.2in}

\lettrine[lines=3]{A} \ STORY HAS ONLY ONE MASTER -- ITS NARR\-ATOR; he decides what he wants his story to do.
I know, I have always known, what I want my stories to achieve -- I want to make people believe.
Believe what I tell.
Believe in it.
Believe me.
Belief is the one effect I'm always looking for, and I apply every device, every pause, every gesture, every verbal nuance and twirl, to that end.
To achieve it, I myself have to believe; if I don't, who will? I must believe ancient Ireland was as I describe it.
The swords really did ring loudly off the shields.
And the armor surely gleamed in the sun.

In fact, I want my listeners to believe so deeply that I almost have them saying to themselves, ``No, he couldn't have been there, that's impossible!'' In order to get them close to that point, I make a great effort to close a specific gap -- the gap that separates the historical fact from the invented tale.
And if they are -- or could be one and the same thing, that's what I call ``the magic of the past,'' and that's what I deal in: the magic of the past.

I know that if my listeners, with their round eyes and rapt faces if they believe that what they're hearing actually happened, they'll remember it forever, especially if the facts are lit in bright lights.
That's why I try and make every tale I tell seem like a film made centuries ago.

The need for belief also accounts for the reason why I tell such different stories.
Some are history, reported by those who were there, and reembroidered by me.
Others are myth, from deep in our souls, handed down by word of mouth, and I am the latest narrator.
And I make up yet others -- because I like the power and the fun of creating worlds.

When it's running well, the effects of a story will last as long as I remain in the presence of my audience.
Then, when they shake themselves back into the present day, and by the time they cross the dark night on their way home, they know and see my sleight of hand.
And I hope they talk and laugh about the evening, as they would have after watching a magician.
I don't even mind if they work out how I did the trick.

\vspace{1em}

\noindent In order to do my job well, to capture their hearts and spirits and minds, I have to know my audience; I have to know a great deal about them.
Every day I learn something new about the Irish and about our island.
Over the years this has built up into a body of knowledge that gives me a good and permanent picture of what we're like as a people.

I know all kinds of things.
For example, I know that we like to fill empty spaces with words.
Sometimes the words don't mean a thing; my father used to murmur every so often, ``Up she flew, and the cock flattened her.'' Out of the blue he'd say it, no context whatever, a meaningless utterance, spoken to no one.
I have no idea why he did it; and I think he merely liked to roll a lively and cryptic phrase in his mouth, to taste it like a bright fruit.

And I also know that, as a people, we generally cohere, that we're a nation of human beings with much in common.
North, south, east, and west, we eat the same foods, feel the same weather, dance the same steps, die of the same diseases.

These common factors make my job easier because I always know to whom I'm speaking.
All through the thirty-two counties of Ireland, most of the adults who have listened to my tales have quit school at the legal age of fourteen.
That is our national pattern since the beginning of the twentieth century.
Even now, I have met very few in the countryside who went to any secondary school, or who then climbed to university level and beyond.
If they have done, it's always been a happenstance of their parents' view and position, rather than a general norm.

For ``ordinary'' people, that is to say the people who work the land and the other people who form the parish, the tradition of education has not yet returned.
The next few decades might restore it to Ireland; it was killed two hundred and fifty years ago by the English in one of their many attempts to eradicate us, and it needs to be brought back to life.

By and large, we have no class structure in Ireland, except among the Protestants.
Many of the Catholic Johnny-come-latelys would like to look down their noses; they'd love to be the haves lording it over the have-nots.
But by dint of cheek and irreverence, and by dint of traditionally being in the same repressed, oppressed boat, Irish Jack is largely as good as his master.

So when I'm addressing such mixed audiences, I have to accommodate a width of response that ranges from illiterate to doctorate.
I live by a guiding principle that I learned in Rome (of which eternal city more another time); ``Never underestimate their intelligence, always underestimate their knowledge.'' And because I know my audience, my little extra details can send shivers through them, especially when I draw on things they know, that they've seen.
If they understand that Irish people two thousand years ago lived similar lives, they grasp my story more powerfully, and they more willingly give themselves over to me.

As you'll have gathered by now, mine is mainly a rural audience.
They think about cattle and meadows, saving hay and reaping corn, tanning hide and spinning wool.
I can be deep in the west, in the snipe-grass smallholdings of poor Mayo, or I can be down south across the richness of Limerick and Tipperary, in what they call the Golden Vale, a seam of rich earth that carves its way across Europe from Hungary.
Or I can be in a house in a town where the church clock chimes.

Wherever I am, my listeners easily recognize from their own lives my descriptions of the ancient Irish -- the farmers who worked the fields, whose wives baked bread for the midday meal, and whose comfortable farmyards housed fat pigs.

\vspace{1em}

\noindent Beneath its broad surface, storytelling should always work hard to say more than it seems to.
When I told you the story of Newgrange, I made the Angry Woman's daughter look like a child with red hair who was there that night.
If I'm describing a beautiful princess.
I'll make sure nobody round that fireside matches her.
The men love the account of the girl, and the women want news of her beautiful gown.

Likewise, before I describe someone ugly.
I'll look carefully around the room in case I come out with an account that offends anyone.
Folk are touchy; it would be easy for some man to feel that I mocked his blemish.

I also do what old storytellers always did.
People need to take a rest now and then, so I make sure that I use little phrases to ease them, to tell them -- and me -- we can relax for a few seconds.
I have composed a wide repertoire.

``And every year, they turned the soil, they sowed their seeds, they welcomed the rain, they hailed the sun, and they were happy.''

``And so, things became plain to everyone.''

``That is one of the ways of the world, and the ways of the world are many.''

``While the ships sailed the seas and the clouds sailed the sky.''

Homer the Greek had such tricks.
When he said, ``Dawn comes early, with rosy fingers,'' and spoke repeatedly of ``the wine-dark sea'' and ``Such were the words that passed between us,'' his audience drew breath with him.
If it was good enough for him twentyseven centuries ago, it's good enough for me.
All of us storytellers, the world over, the last few of us still traveling -- we play with such chosen phrases.
They allow us to change gear, and we can also use them to tighten the screw as the story begins to climb.

\vspace{1em}

\noindent In some ways I'm very like an actor; this is a performance that I give, I can think of no other word for it.
It takes place in a very intimate theater, with no stage and no proscenium arch and no curtains and an audience that is often only inches away.
To address those circumstances, much of what I do has to be theatrical, and every actor in the world needs his props.

My main prop is my pipe.
First, I fill it with tobacco as carefully as a man counting his money.
Next, I strike the match -- I love how that little yellow flash surges.
Then I suck the flame down into the bowl until the pipe has swallowed its light.
I like to think it gives the impression that I'm inhaling the sorcery of fire.
And I tamp down the hot tobacco with my fingertip; this looks fearless, and I want people to think.
Here's a man who can touch flames and never burn.
Not a word do I say through all of this, and slowly the listeners fall silent as they wait for me to speak.

When finally I look up from the pipe, I make my eyes rove through the audience, because I want to be sure that every person is ready to come happily into my grasp.
A storyteller has to know how to get power over people.

I also use my hat as a prop, my old hat that keeps out the rain and deflects the sun, on whose brim the dew settles like little pearls on a necklace.
Once it was a homburg, but it has been long bruised by the world.

Still, it is my good theatrical friend.
I take it off when I want to build a pause or unwrap a detail.
Sometimes I pass the brim through my hands as if attempting to shape it into a consistent straight line.
Or I fiddle with the decrepit old ribbon, which slides easily around the crown, like a belt on a man grown thin.
Or I punch it gently until the homburg shape becomes one large dome, like a comical hat.

And then with the heel of my hand I restore the soft dent in the crown and return the hat to my head.
More than one child has said to me that they think the hat somehow shelters my tales, whose colors would otherwise come bursting out of my cranium like the rays of rainbows.

At the end of the story I follow a long-practiced routine.
I lift my hat from my head and replace it, as a gentleman would to a lady; I take my dead pipe from my mouth; I lean down and tap it on the wide hearthstone.
When the small gray powder falls out -- I always smoke my pipe down to the last ash -- I like to think it looks as though I wrested every last fiber from the tale.

\vspace{1em}

\noindent Then there is the language of the body to aid the delivery of the words.
Musicians know how and when to arrange the spaces between their notes.
So should storytellers.
I pause and let the audience watch me, and I wait for their minds to come to me.
Or I turn my head away from the room and gaze into the fire.
Or I look around at the assembled listeners, almost as if seeking someone.
And then I stare once more into the flames.
When at last I turn to face the room fully, I know that my face seems to glow red and my eyes black with the fire I have absorbed.

All I'm saying is, I use the techniques of an actor, with my body as my tool bench and my hands alternating between my pipe, my hat, and my gestures of expression.
And I have one special gesture, which I use almost every night.

As my arms open in wide movements, I make my eyes flash like a mesmerist.
While my left hand grips the pipe bowl tight as desire, my right hand closes in a fist.
Slowly, slowly, like the petals of some bony flower, my gnarled fingers begin to open, so that the hand becomes like a little bowl -- and then the fingers flatten and extend to their lengths until the hand itself lies flat, upward and open in the air.
I always make this gesture at the moment when I feel the audience most hushed.

At that moment I would not be surprised to find sitting on my hand the population of the room, silent and enthralled -- and all permanently half-smiling at me and my magical words.
I want them there, I want them like that, I want them to say, ``He held his audience in the palm of his hand'' -- and all my efforts, my gestures and expressions, are directed to that aim.
When they are there -- they believe.

\end{document}